%------------------------------------------------------------------------------
\subsection{\file{models/kazmierczak2024/water\_flux.jl}: from the potential to $q$ and $W$}\label{ref:waterflux}
\emph{Tour PART 14.}

The core of K24: the potential, filling, gradients, smoothing, the four routing traversals, the Picard loops and the
thickness closures. Its 1300 lines are mostly docstrings recording why each step is done the way it is.

\subsubsection*{The driver and the source terms}
\begin{fnlist}
\item[update\_q!(model, grid, state)] \code{update\_phi0!}; \code{prepare\_routing!} for the scheme (\ref{ref:routing}); \code{corfac}
  (Eq.~\ref{eq:corfac}, with \code{eps(T)} in the denominator, computed as a scalar outside the broadcast so a Float32 model is not
  promoted, and \code{x*x} rather than \code{x\^{}2.0}); \code{invalidate\_routing\_tape!} (the gradients just changed); then
  \code{resolve\_q!(model, grid, state, model.dissipation\_melt, model.sliding\_law)}. The docstring still says $q$ is clamped at
  $10^5$\,m$^2$/yr; that is now opt-in (\code{q\_max}, default $\infty$).
\item[add\_dissipation\_term!(model, ::DissipationMeltOff / ::DissipationMeltOn)] Off: \code{Q\_diss .= 0}. On: forwards to the
  discretisation's method in \file{routing.jl}.
\item[add\_friction\_term!(model), add\_friction\_term!(model, ::CellCentredFriction)] Forwards to the model's friction discretisation;
  the cell-centred method sets $Q_b = \tau_b|u_b|$ and adds $Q_b/L_w$ to \code{mdot\_total}. \code{model.tau\_b} is always current when it
  is called (every \code{resolve\_q!} updates it first).
\end{fnlist}

\subsubsection*{Routing traversals}
\begin{fnlist}
\item[route\_psi\_out!(model, grid, state[, algorithm])] Dispatches on \code{model.psi\_out\_algorithm} to \code{update\_psi\_out!},
  \code{update\_psi\_out\_iterative!}, \code{update\_psi\_out\_topological!} or \code{update\_psi\_out\_taped!}. Every \code{resolve\_q!} calls
  this, never a traversal directly.
\item[invalidate\_routing\_tape!(model)] \code{routing\_tape.valid = false}. Call it before \code{route\_psi\_out!} if you changed the mask or
  the gradients yourself.
\item[update\_psi\_out\_taped!(model, grid, state)] Records the tape if it is not valid, then replays it with the current
  \code{mdot\_total}.
\item[record\_routing\_tape!(tape, model, grid, state)] Calls \code{record\_routing\_tape\_kernel!} (GDSWarner) or
  \code{record\_routing\_tape\_weights\_kernel!} (\file{routing.jl}, other schemes), warns once if the cap was hit (from here, not inside
  the kernel: the source measured that a logging macro in the hot loop made the traversal 6 times slower), and sets
  \code{valid}.
\item[push\_tape\_op!(tape, ci, cj, ni, nj, w)] Appends one operation.
\item[record\_routing\_tape\_kernel!(tape, mask, sx, sy, Nx, Ny, dx, dy, epsT, max\_calls)] The explicit-stack traversal of
  \code{update\_psi\_out\_iterative!}, recording the operations instead of performing them. Frames $(i,j,k)$: $k = 0$ not yet visited,
  $1\le k\le4$ next neighbour, $k = 5$ done (push the clamp). Neighbours without grounded ice are skipped, not recorded as
  \code{+= 0}. Returns whether the cap was hit.
\item[replay\_routing\_tape!(psi\_out, mdot\_total, mask, tape, Nx, Ny, dx, dy, rho\_w)] Sets every grounded cell to its own source
  $\mdot\Dx\Dy/\rhow$ (the same expression and evaluation order as the recursion), then applies the operations in order. A separate
  function taking the arrays so the loop specialises on their types. A cell's \code{psi\_out} is never read before its first visit, so
  setting all sources up front changes no value read.
\item[update\_q\_from\_psi\_out!(model)] \code{QFromOutflow}: $q = \min(\max(\psiout/(c_f+\varepsilon), q_{\min}), q_{\max})$. The
  three-argument method is in \file{routing.jl}.
\item[routing\_weight(sx, sy, di, dj, dx, dy, epsT)] GDSWarner's fraction toward neighbour $(d_i, d_j)$:
  $(s_x d_i\Dy + s_y d_j\Dx)/(|s_x|\Dy + |s_y|\Dx + \varepsilon)$, positive only toward the flow. The positive fractions sum to 1.
\item[accumulate\_psi\_out!(model, i, j, grid, state, call\_count)] The recursion (Le Brocq's and KORI-ULB's algorithm): a non-grounded
  cell returns 0; a visited cell returns its current \code{psi\_out}; otherwise mark visited, set the source, count the visit (past
  \code{max\_psi\_out\_calls}: warn once, clamp, return), add \code{accumulate\_psi\_out!(neighbour) * w} for every neighbour with
  $w>0$, clamp at 0. Returns \code{zero(eltype(...))} rather than \code{0.0}: the source records, from \code{@code\_warntype}, that a
  \code{Float64} literal made the return type a \code{Union} for a Float32 model.
\item[update\_psi\_out!(model, grid, state)] Clears \code{visited} and starts the recursion from every grounded cell, with one
  \code{Ref} counter for the whole sweep.
\item[update\_psi\_out\_iterative!(model, grid, state)] The same traversal with an explicit stack (a \code{Vector}), bounded only by
  heap memory. When an unvisited neighbour is found it is pushed without advancing the parent's $k$, so on return the parent
  revisits the same neighbour, now visited, and folds it in: exactly how a return value reaches a paused caller.
\item[update\_psi\_out\_topological!(model, grid, state, allow\_cycles)] Builds each grounded cell's outgoing edges and every cell's
  in-degree in one pass, then Kahn's algorithm with an array-backed queue (a moving head index, since \code{popfirst!} on a
  \code{Vector} is $O(n)$). Cells inside a cycle never reach in-degree 0: then an error, or with \code{allow\_cycles} a warning and
  those cells left with only their partial upstream inflow (their own source never added). Needs no visit cap.
\end{fnlist}

\subsubsection*{The fixed points}
\begin{fnlist}
\item[resolve\_q!(model, grid, state, ::DissipationMeltOff, ::Union\{NoFriction, PrescribedField, Weertman\})] One pass:
  $\tau_b$, \code{mdot\_total = mdot\_fixed + i\_eb}, friction term, route, $q$.
\item[resolve\_q!(..., ::DissipationMeltOn, ::Union\{NoFriction, PrescribedField, Weertman\})] Picard on the dissipation: $\tau_b$ once;
  then each sweep rebuilds the source from the current $q$, routes, converts, and stops when
  $\max|q - q_{\rm prev}| \le$ \code{dissipation\_rtol}$\cdot\max|q|$ over grounded cells (default $10^{-12}$), or after
  \code{max\_dissipation\_iters} = 20 sweeps. \code{dissipation\_verbose} logs the outcome. $q$ carries over from the previous call
  (zero on the first), so the first sweep has no dissipation.
\item[resolve\_q!(..., ::AbstractDissipationMelt, ::AbstractPressureDependentSlidingLaw)] The joint loop: each sweep recomputes $\tau_b$
  from the current $N$, adds friction and (if on) dissipation, routes, converts, then calls \code{update\_N!}; stops when both $q$ and
  $N$ change by less than \code{qN\_rtol} (default $10^{-8}$) relative to their maxima, or after \code{max\_qN\_iters} = 20 sweeps.
  $N$ starts from \code{state.N} (zero on a fresh state, so the first sweep has no friction). Section~\ref{sec:fixedpoints}.
\end{fnlist}

\subsubsection*{The water thickness}
\begin{fnlist}
\item[update\_W!(model, grid, state[, algorithm])] Dispatches on \code{model.water\_thickness\_algorithm}. The
  \code{ArealConduitThickness} method writes $\Sinf/l_c$ unclamped (the bounds were chosen for a thin sheet, and a dense conduit
  network smeared over its spacing can exceed them).
\item[update\_W\_darcy\_weisbach!(model, grid, state, ::LocalGradient / ::MeanGradient)] $(f\rhow q^2/(4|\nabla\phi_0| + \varepsilon))^{1/3}$,
  clamped; \code{q*q} rather than \code{q\^{}2} (libm \code{pow}, and \code{@.}'s \code{literal\_pow} rewrite).
\item[update\_W\_laminar!(model, grid, state, ::MeanGradient / ::LocalGradient)] $(12\eta_w q/|\nabla\phi_0|_s)^{1/3}$, clamped; the mean
  variant has no $\varepsilon$ guard.
\end{fnlist}

\subsubsection*{Potential, filling, gradients, smoothing}
\begin{fnlist}
\item[update\_phi0!(model, grid, state)] Eq.~\eqref{eq:phi0}.
\item[potential\_filling!(model, grid, state; from\_true\_potential = true)] Copies $\phi_0$ into \code{phi0\_filled} (unless the caller
  already put the surface there: the smoothed potential of the non-GDS schemes), fills it with \code{fill\_potential!}, and sets
  \code{model.h} $= (\phif - \rhow g b)/(\rhoi g)$.
\item[fill\_potential!(model, grid, state, ::Union\{JacobiFill, LowestNeighbourFill\})] The candidate-list Jacobi passes
  (Section~\ref{disc:filling}): a first full-grid pass writes the new values of strict minima to \code{phi0\_tmp} and lists them; each
  later pass copies them back, collects the filled cells and their neighbours, and re-tests only those; stops when a pass fills
  nothing or after \code{fill\_iters} passes. Same result as full-grid passes, bit for bit (tour).
\item[fill\_potential!(model, grid, state, ::PriorityFloodFill)] Section~\ref{disc:filling}. Reuses \code{fill\_stamp} as the ``queued''
  flag; the heap is two parallel vectors with a size hint of a quarter of the grid.
\item[heap\_push!(p, k, pv, kv), heap\_pop!(p, k)] A minimal binary min-heap on (priority, linear index).
\item[fill\_if\_strict\_minimum!(phi0, phi0\_tmp, filled, i, j, Nx, Ny, lowest = false)] If $(i,j)$ is lower than all four (edge-replicated)
  neighbours, writes their mean (or minimum) to \code{phi0\_tmp} and records the cell.
\item[collect\_fill\_candidates!(candidates, stamp, pass, filled, Nx, Ny)] The filled cells and their in-domain neighbours as
  duplicate-free linear indices, deduplicated by stamping each with the pass number.
\item[update\_potential\_gradients!(model, grid)] \code{abs\_grad\_phi0} from the true $\phi_0$ (Euclidean), then
  \code{minus\_grad\_phi0\_x/y} from $\phif$ (the buffers are reused).
\item[coupling\_kernel(model, grid, h\_avg)] The normalised cone of Section~\ref{disc:smoothing}. The branch for a cone narrower than
  the finer spacing changes \code{scale} but not \code{width}, so the kernel is $1\times1$ (no smoothing) whatever it does
  (Section~\ref{pit:open}).
\item[update\_smoothed\_potential\_gradients!(model, grid, state)] GDS schemes: convolves both gradient components with the cone for
  $\bar h = \max(\text{masked mean of } \code{model.h}, 10)$; sets \code{abs\_grad\_phi0\_s} $= |s_x| + |s_y|$. With
  \code{longcoupwater == 0} copies the raw components.
\end{fnlist}

%------------------------------------------------------------------------------
\subsection{\file{models/kazmierczak2024/routing.jl}: routing schemes, face fluxes, freeze-on}\label{ref:routing}
\emph{Tour PART 15.}

The Le Brocq et al.\ (2006) schemes as outflow fractions, the face fluxes and what is built from them, and the freeze-on
capacity. The header comment still calls \code{GDSWarner} the default; since the change to Warner it is the original
scheme kept for reproduction.
\begin{fnlist}
\item[ROUTE\_OFFSETS, ROUTE\_OPPOSITE, TARBOTON\_FACETS] The 8 directions (W, E, S, N, SW, SE, NW, NE), their reverses, and the 8
  Tarboton facets as (cardinal, adjacent diagonal) pairs.
\item[in\_domain(i, j, Nx, Ny)] Bounds test.
\item[needs\_face\_fluxes(model)] \code{QFromFaceAverage}, or \code{FaceDissipation} with the dissipation on.
\item[prepare\_routing!(model, grid, state, ::Union\{GDSWarner, GDSTarboton\})] Fill, gradients, smoothed gradients; the outflow
  fractions only for GDSTarboton or when face fluxes are needed.
\item[prepare\_routing!(model, grid, state, ::Union\{Warner, Quinn, Tarboton, ModifiedTarboton\})] Convolve $\phi_0$ with the cone (for
  $\bar h$ from \code{state.h}, not the filled thickness) into \code{phi0\_filled}, fill it in place, take gradients, copy them to
  \code{minus\_grad\_phi0\_sx/sy} (used for $c_f$ and by ModifiedTarboton), compute the fractions.
\item[compute\_routing\_weights!(model, grid, state), routing\_weights\_kernel!(W, scheme, phi, sx, sy, mask, Nx, Ny, dx, dy)] Zero
  \code{w8}, then \code{cell\_weights!} for every grounded cell.
\item[cell\_weights!(W, scheme, i, j, phi, sx, sy, Nx, Ny, dx, dy)] One method per scheme, Section~\ref{disc:schemes}. Fractions sum to
  1, or to 0 at a cell with nowhere downhill to go (a sink, which keeps its water).
\item[angle\_split!(W, i, j, fx, fy, dx, dy)] Splits one direction $\theta$ between the two of the 8 neighbours whose angles
  ($0$, $\arctan(\Dy/\Dx)$, $\pi/2$, \ldots) bracket it, linearly in angle.
\item[record\_routing\_tape\_weights\_kernel!(tape, mask, W, ndirs, Nx, Ny, max\_calls)] The generic recorder: cell $c$ receives
  \code{psi\_out[n] * W[opposite(d), n]} from neighbour $n$ in direction $d$, over the first \code{ndirs} directions.
\item[update\_face\_fluxes!(model, grid, state), face\_fluxes\_kernel!(Fx, Fy, diss, psi, W, mask, phi, Nx, Ny, dx, dy, L\_w, do\_diss)] The
  net face fluxes of Section~\ref{disc:qconv}, and with \code{FaceDissipation} the cell dissipation melt of
  Section~\ref{disc:dissipation}.
\item[update\_q\_from\_psi\_out!(model, grid, state)] Updates the face fluxes if needed, then \code{q\_from\_psi\_out!}.
\item[q\_from\_psi\_out!(model, grid, state, ::QFromOutflow / ::QFromFaceAverage), q\_from\_faces\_kernel!] The two conversions; the face
  one writes 0 off the grounded mask before clamping.
\item[add\_dissipation\_term!(model, ::CellCentredDissipation / ::FaceDissipation)] $Q_{\rm diss} = |q\,|\nabla\phi_0||$, or
  $Q_{\rm diss} = L_w\cdot$\code{tape.diss}; both add the melt to \code{mdot\_total}. The face version uses the face fluxes of the
  previous sweep, so its value lags by one sweep and is exact at convergence.
\item[freeze\_on\_capacity!(C, model, grid, state), freeze\_on\_capacity\_kernel!] Section~\ref{disc:freeze}. For GDSWarner without
  face fluxes it first fills \code{w8}. Returns \code{C}.
\end{fnlist}
